% ══════════════════════════════════════════════════════════════════════════════
%  EXPERIMENT 2 — HOUSEHOLD ENERGY AUDIT                                (Lesson 6)
%  Classroom survey in the lesson; the home audit is homework.
% ══════════════════════════════════════════════════════════════════════════════
\experimentdiv{Household Energy Audit}{Power ratings \textbullet{} Energy in kWh \textbullet{} Season and building design \textbullet{} A recommendation}{6}\label{sec:2}
\lessoncodes{VC2S8U16 \textbullet{} VC2S8H04 \textbullet{} VC2S8I02, I04, I05, I07, I08}

\vocabhead
\begin{ctab}
\begin{tabular}{|SW{4.2cm}|W{11.6cm}|}
  \hline\hd & \thead{Word} & \thead{Definition (copy from the slide)} \\ \hline
  & \scaf{\textbf{Energy audit}} & \scaf{A survey of where energy is used, how much, and how it could be reduced.} \\ \hline
  & \textbf{Power} & \\ \hline
  & \textbf{Kilowatt-hour} & \\ \hline
  & \textbf{Energy rating label} & \\ \hline
  & \textbf{Standby} & \\ \hline
\end{tabular}
\end{ctab}

\begin{lipanel}
\inquiryli{Collect data in a table, find the biggest energy users, and make a recommendation.}
\understandli{Explain how energy use depends on appliances, building design, season and climate.}
\sctitle
\begin{itemize}[leftmargin=1.4em, itemsep=0pt, topsep=1pt]
  \item I can calculate the energy an appliance uses in a day, in kWh.
  \item I can identify the biggest energy users in a home and explain why.
  \item I can write a recommendation to a household, backed by my data.
\end{itemize}
\end{lipanel}

\begin{skillsfocus}
\focus{PC5}{Audit a home with permission and privacy.}
\focus{PA1}{Organise appliance data in a table.}
\focus{PA2}{Draw a column graph of energy use.}
\focus{PA3}{Calculate energy $=$ power $\times$ time.}
\focus{PA6}{Does colder weather cause more use?}
\focus{EV4}{Judge energy advice and labels.}
\focus{CO1}{Write a recommendation backed by data.}
\end{skillsfocus}

% ─── TECHNIQUE ────────────────────────────────────────────────────────────────
\exppart{Technique spotlight}{Blue}

\begin{technique}{Reading power and energy labels}{PC6}
\begin{minipage}[c]{0.6\linewidth}
\textbf{Power} is how fast an appliance uses energy. It is printed on the appliance's rating plate, in watts (W) or kilowatts (kW): 1\,kW $=$ 1000\,W.

\smallskip
Many appliances also carry an \textbf{energy rating label}. More stars means the appliance uses less energy than others of its size. The number on the label is its estimated energy use in a year, in kilowatt-hours (kWh).
\end{minipage}\hfill
\begin{minipage}[c]{0.36\linewidth}\centering
\begin{tikzpicture}[scale=0.85]
  \draw[eq, rounded corners=2mm] (0,0) rectangle (4,4.2);
  \node[font=\headingfont\bfseries\footnotesize] at (2,3.85) {ENERGY RATING};
  \foreach \k in {0,...,5}{
    \pgfmathsetmacro\ang{160-\k*28}
    \draw[Ink, line width=0.8pt, fill=\ifnum\k<4 Ink\else white\fi] ($(2,1.6)+(\ang:1.6)+(90:0.2)$) -- ($(2,1.6)+(\ang:1.6)+(162:0.2)$) -- ($(2,1.6)+(\ang:1.6)+(234:0.2)$) -- ($(2,1.6)+(\ang:1.6)+(306:0.2)$) -- ($(2,1.6)+(\ang:1.6)+(18:0.2)$) -- cycle;}
  \node[font=\headingfont\bfseries\Large] at (2,1.25) {350};
  \node[font=\scriptsize] at (2,0.75) {kWh per year};
  \node[font=\scriptsize] at (2,0.35) {4 stars out of 6};
\end{tikzpicture}
\end{minipage}
\end{technique}

\begin{technique}{Calculating the energy an appliance uses}{PA3}
\textbf{Energy (kWh)} $=$ \textbf{power (kW)} $\times$ \textbf{time (hours)}. \quad Change watts to kilowatts by dividing by 1000, and minutes to hours by dividing by 60.
\end{technique}

\begin{workedex}{A kettle}
A 2000\,W kettle is used for 15 minutes a day. How much energy does it use in a day?

\smallskip
\textbf{Power:} $2000 \div 1000 = 2$\,kW. \quad \textbf{Time:} $15 \div 60 = 0.25$\,h.\\
\textbf{Energy:} $2 \times 0.25 = $ \textbf{0.5\,kWh} per day.
\end{workedex}

\begin{wedo}{A television}
A 120\,W television is on for 4 hours a day. How much energy does it use in a day?
\workbox{3.5cm}
\end{wedo}

\begin{youdobox}[Energy per day]
\task Complete the table. Then find the cost of each appliance per day, if electricity costs 30 cents per kWh. \skr{PA3}
\begin{ctab}
\begin{tabular}{|SW{3.4cm}|Z{2.2cm}|Z{2.4cm}|Z{3.6cm}|Z{3.4cm}|}
  \hline\hd & \thead{Appliance} & \thead{Power (W)} & \thead{Time per day (h)} & \thead{Energy per day (kWh)} & \thead{Cost per day (cents)} \\ \hline
  & \scaf{e.g. kettle} & \scaf{2000} & \scaf{0.25} & \scaf{$2 \times 0.25 = 0.5$} & \scaf{$0.5 \times 30 = 15$} \\ \hline
  & LED light globe & 10 & 5 & & \\ \hline
  & Electric heater & 2000 & 3 & & \\ \hline
  & Phone charger & 5 & 3 & & \\ \hline
\end{tabular}
\end{ctab}
\end{youdobox}

\begin{technique}{Judging energy advice}{EV4}
Energy advice is everywhere: in advertisements, on websites, from friends. Before you trust it, ask: \textbf{Who} wrote it? \textbf{Why} --- are they selling something? \textbf{What evidence} does it give? \textbf{When} was it written? Government energy rating labels are tested in the same way for every appliance, so they let you compare fairly.
\end{technique}

\begin{workedex}{Two sources}
\emph{Source 1:} an advertisement says ``Our heater is the cheapest to run!'' \emph{Source 2:} a government website lists the yearly energy use of every heater, tested the same way.

\smallskip
Source 2 is more reliable: it has no product to sell, and every heater is tested by the same method, so the numbers can be compared.
\end{workedex}

\begin{wedo}{Two sources}
A friend says that leaving the TV on standby uses no electricity at all. A government energy website says standby can make up about 3\% of a home's electricity use. Which would you trust, and why?
\answerbox{4.5cm}
\end{wedo}

% ─── PLAN AND SAFETY / METHOD ────────────────────────────────────────────────
\exppart{Plan, safety and method}{Teal}

\textbf{Aim:} To find out where energy is used in our classroom and in a home, and to recommend ways to use less.

\begin{infobox}{Ethics and safety}
\textbf{Ethics:} only audit your own home, and only with permission from the adults who live there. Do not write down names, addresses or bills in this booklet. \skr{PC5}\par
\textbf{Safety:} only read labels you can see easily. Do not unplug, move or open any appliance, and never touch a switch or plug with wet hands. Ask an adult if a label is hidden.
\end{infobox}

\begin{infobox}{Typical power ratings --- use these if you cannot read a label}
\renewcommand{\arraystretch}{1.15}\small
\begin{tabular}{@{}l r @{\hspace{8mm}} l r @{\hspace{8mm}} l r@{}}
Fridge (average) & 150\,W & Oven & 2400\,W & LED globe & 10\,W \\
Television & 120\,W & Washing machine & 500\,W & Halogen globe & 50\,W \\
Kettle & 2000\,W & Clothes dryer & 2400\,W & Laptop & 40\,W \\
Toaster & 900\,W & Electric heater & 2000\,W & Desktop computer & 100\,W \\
Microwave & 1000\,W & Air conditioner & 2500\,W & Game console & 150\,W \\
\end{tabular}
\end{infobox}

\subsubsection*{Method --- Part 1: The classroom (in class)}
\begin{enumerate}[itemsep=2pt]
  \item List every appliance that uses electricity in the classroom, and how many there are.
  \item Find the power of each, from its label or the table of typical ratings.
  \item Estimate how many hours a day each is switched on. Record everything in Table 2.1.
\end{enumerate}

\subsubsection*{Method --- Part 2: Your home (homework)}
\begin{enumerate}[itemsep=2pt]
  \item With an adult, list ten appliances in your home and find the power of each.
  \item Ask how many hours a day each is used in winter and in summer. Record these in Table 2.2.
  \item Complete the building design checklist (Table 2.3) together.
  \item Back in class, calculate the energy each appliance uses per day.
\end{enumerate}

% ─── RESULTS ─────────────────────────────────────────────────────────────────
\exppart{Results}{Amber}

\subsubsection*{Table 2.1 --- The classroom \skr{PA1}}
\begin{ctab}
\begin{tabular}{|SW{4.4cm}|Z{1.6cm}|Z{2.6cm}|Z{3.0cm}|Z{3.4cm}|}
  \hline\hd & \thead{Appliance} & \thead{Number} & \thead{Power of each (W)} & \thead{Hours on per day} & \thead{Energy per day (kWh)} \\ \hline
  & \scaf{e.g. LED ceiling light} & \scaf{12} & \scaf{20} & \scaf{6} & \scaf{$12 \times 0.02 \times 6 = 1.44$} \\ \hline
  & & & & & \\ \hline
  & & & & & \\ \hline
  & & & & & \\ \hline
  & & & & & \\ \hline
\end{tabular}
\end{ctab}

\clearpage
\subsubsection*{Table 2.2 --- My home}
\begin{ctab}
\begin{tabular}{|SW{3.6cm}|Z{2.0cm}|Z{2.2cm}|Z{2.6cm}|Z{2.2cm}|Z{2.6cm}|}
  \hline\hd & \thead{Appliance} & \thead{Power (W)} & \thead{Winter hours per day} & \thead{Winter kWh per day} & \thead{Summer hours per day} & \thead{Summer kWh per day} \\ \hline
  & \scaf{e.g. heater} & \scaf{2000} & \scaf{4} & \scaf{8} & \scaf{0} & \scaf{0} \\ \hline
  & & & & & & \\ \hline
  & & & & & & \\ \hline
  & & & & & & \\ \hline
  & & & & & & \\ \hline
  & & & & & & \\ \hline
  & & & & & & \\ \hline
  & & & & & & \\ \hline
  & & & & & & \\ \hline
  & & & & & & \\ \hline
  & & & & & & \\ \hline
\end{tabular}
\end{ctab}

\clearpage
\subsubsection*{Table 2.3 --- Building design checklist}
\begin{ctab}
\begin{tabular}{|SW{6.0cm}|Z{2.2cm}|W{6.8cm}|}
  \hline\hd & \thead{Feature} & \thead{Yes / No / Not sure} & \thead{How it changes energy use} \\ \hline
  & \scaf{Curtains or blinds on the windows} & \scaf{Yes} & \scaf{Keep heat in on winter nights and out on summer days.} \\ \hline
  & Ceiling insulation & & \\ \hline
  & Draught seals on doors & & \\ \hline
  & Double-glazed windows & & \\ \hline
  & Main living area has north-facing windows & & \\ \hline
  & All light globes are LEDs & & \\ \hline
  & Solar panels on the roof & & \\ \hline
\end{tabular}
\end{ctab}

\needspace{16\baselineskip}
\task Draw a column graph of the winter energy use per day of your six biggest users from Table 2.2. Label both axes and give the graph a title. \skr{PA2}
\graphgrid{9cm}

% ─── ANALYSIS ────────────────────────────────────────────────────────────────
\exppart{Analysis and evaluation}{Coral}

\task Which three appliances in your home use the most energy in winter? Suggest why they use so much. \skr{PA5}
\answerbox{5cm}

\needspace{14\baselineskip}
\task Compare your home's energy use in winter and summer. \skr{PA6}
\pt{a}Total energy per day in winter and in summer:
\workbox{3cm}
\pt{b}Energy use goes up as the weather gets colder. Is the cold weather itself the cause? Explain what actually uses the energy.
\answerbox{5cm}

\needspace{10\baselineskip}
\task Melbourne has cold winters and hot summers. Explain how a home in Darwin, where it is hot all year, would use energy differently.
\answerbox{5cm}

\needspace{10\baselineskip}
\task Choose two features from Table 2.3. Explain how each one changes the energy a home uses for heating and cooling.
\answerbox{5cm}

\needspace{12\baselineskip}
\task A family is choosing between two fridges of the same size. Fridge A has 3 stars and uses 450\,kWh a year. Fridge B has 5 stars and uses 300\,kWh a year. \skr{EV4}
\pt{a}How much energy would Fridge B save in a year?
\workbox{2.5cm}
\pt{b}Explain how the energy rating label helps the family choose.
\answerbox{5cm}

\needspace{12\baselineskip}
\task Write a recommendation to your household: three changes that would reduce its energy use. Back each one with data from your audit, and write in the third person. \skr{CO1}

\scaf{The audit found that the household uses the most energy for \dots. It is recommended that \dots, because \dots}
\answerbox{8cm}

\needspace{10\baselineskip}
\task Write a new question about household energy use that you could investigate next. \skr{EV5}
\answerbox{3cm}

\begin{yr8box}
\begin{keepext}
\exttask A household replaces eight 50\,W halogen globes with 10\,W LED globes. The lights are on for 5 hours a day. Calculate the energy saved in a year, and the money saved at 30 cents per kWh.
\workbox{4cm}
\end{keepext}
\begin{keepext}
\exttask Solar panels make the most electricity in the middle of the day, but many homes use the most electricity in the evening. Suggest two ways a household could make better use of its solar panels.
\answerbox{5cm}
\end{keepext}
\end{yr8box}

% ─── SKILL CHECK ─────────────────────────────────────────────────────────────
\exppart{Skill check}{NavyBlue}

\begin{skillcheck}
\scrow{PC5}{I audited a home with permission, and kept private information out of my booklet.}
\scrow{PA1}{I organised my appliance data in tables with units in the headings.}
\scrow{PA2}{I drew a column graph with labelled axes and a title.}
\scrow{PA3}{I can calculate energy with power $\times$ time.}
\scrow{PA6}{I can tell whether one thing causes another, or only changes with it.}
\scrow{EV4}{I can judge whether energy advice can be trusted.}
\scrow{CO1}{I can write a recommendation backed by data.}
\end{skillcheck}

\task Fill in the \emph{End of unit} column of your skills tracker at the front of this booklet. Which skill has improved most in this unit, and which one will you work on next?
\answerbox{5cm}

\end{document}
